% SCHO research specification; experimental validation is pending.
% This is a pure LaTeX chapter without a document preamble.




\section{Motivation: Function Coordinates Instead of Fixed Sensors}
\label{sec:scho-motivation}

We propose \emph{Spline-Coefficient HyperOperator} (SCHO): represent both
the input and output in fixed finite-dimensional spline spaces, then
learn an operator between the coefficient vectors. This is different
from a fixed point-sampling encoder followed by a query-coordinate
decoder. It is a finite-dimensional model, not a representation of
all functions in an infinite-dimensional space.

Fix domains $X,Y$ and linearly independent B-spline families
$\{B_i\}_{i=1}^{m}$ on $X$ and $\{\Phi_j\}_{j=1}^{n}$ on $Y$.
The knot vectors, degrees, boundaries, and basis order must be specified
before training. Define
\begin{equation}
  V_m=\operatorname{span}\{B_1,\ldots,B_m\},\quad
  U_n=\operatorname{span}\{\Phi_1,\ldots,\Phi_n\}.
  \label{eq:scho-spaces}
\end{equation}
For the input function and learned output, set
\begin{equation}
  f_\theta(x)=\sum_{i=1}^{m}\theta_i B_i(x),\quad
  \beta=H_\alpha(\theta),\quad
  u_\beta(y)=\sum_{j=1}^{n}\beta_j\Phi_j(y).
  \label{eq:scho-base}
\end{equation}
The hypernetwork $H_\alpha:\mathbb R^m\to\mathbb R^n$ can be a small
ordinary MLP. The output is a whole spline function; its value at
any query coordinate is obtained through the fixed basis evaluation.
Only the use of B-spline expansion is borrowed from one component of
Kolmogorov--Arnold Networks (KANs), not the full KAN architecture.




\section{Coordinate Uniqueness and Its Limits}
\label{sec:scho-uniqueness}

For fixed, linearly independent bases,
\begin{equation}
  f_{\theta_1}=f_{\theta_2}\iff\theta_1=\theta_2,\qquad
  u_{\beta_1}=u_{\beta_2}\iff\beta_1=\beta_2.
  \label{eq:scho-unique}
\end{equation}
This is uniqueness of \emph{coordinates in the selected spline spaces}.
Representing arbitrary functions outside those spaces requires a
specified projection, interpolation, or fitting procedure, with its own
approximation and conditioning errors. These statements do not imply
uniqueness of the neural network weights $\alpha$ or PDE solutions.

In a simple integration example, when the output basis spans the exact
antiderivatives of the input basis,
\begin{equation}
  \mathcal I f_\theta(x)=\int_0^x f_\theta(t)\,dt
  =\sum_{j=1}^{n}(A\theta)_j\Phi_j(x).
  \label{eq:scho-integral}
\end{equation}
The exact coefficient map $A$ is then a fixed matrix, so learning it is
a useful \emph{correctness check}, not evidence that a hypernetwork
accelerates integration.




\section{SCHO-Res: A Restricted Correction Space}
\label{sec:scho-res}

To enrich the main spline space, choose an ambient output Hilbert space
$V\supset U_n$ with a documented inner product and let
$P_n:V\to U_n$ be its orthogonal projection. For chosen candidate
functions $\widetilde\psi_k\in V$, define
\begin{equation}
  \psi_k=(I-P_n)\widetilde\psi_k.
  \label{eq:scho-projection}
\end{equation}
Discard zero or dependent projected features. With
$W_r=\operatorname{span}\{\psi_1,\ldots,\psi_r\}$,
we have $W_r\subset U_n^\perp$ and $U_n\cap W_r=\{0\}$.
A small Dense network produces correction coefficients
$D_\varphi:\mathbb R^m\to\mathbb R^r$:
\begin{equation}
  \boxed{
    \begin{aligned}
      \mathcal H_{\alpha,\varphi}:\theta
       & \longmapsto\bigl(H_\alpha(\theta),D_\varphi(\theta)\bigr)
      =(\beta,\gamma),                                             \\
      \widehat u_\theta(y)
       & =\sum_j\beta_j\Phi_j(y)+\sum_k\gamma_k\psi_k(y).
    \end{aligned}}
  \label{eq:scho-res-operator}
\end{equation}
Because the two bases together are linearly independent, the output
function has a unique $(\beta,\gamma)$ representation. This does not
fix the internal neural-network weight symmetries.

Alternatively a Dense neural field $N_\varphi(\theta,\cdot)$ can be
projected as $(I-P_n)N_\varphi(\theta,\cdot)$. This guarantees only
function-level complementarity, not unique network parameters.
The fixed correction-basis version is the first experiment.

For homogeneous Dirichlet data, use trial and correction functions in
$H_0^1(0,1)$, where the energy inner product
$a(v,w)=\int_0^1v'w'\,dx$ is valid. The corresponding projection is
computed via a fixed Gram matrix. For inhomogeneous boundary data,
a boundary lift is required. Exact function-space orthogonality can
also be affected by numerical projection and quadrature errors.





\subsection{A worked example: degree-one splines and exact integration}
\label{sec:scho-concrete-integration}

Start with any partition of $[0,1]$:
\begin{equation}
  0=x_0<x_1<\cdots<x_N=1,\qquad h_i=x_{i+1}-x_i.
\end{equation}
Degree-one clamped B-splines $B_i$ are exactly the nodal
piecewise-linear hat functions used in one-dimensional $P_1$
finite elements. Their knot vector is
$(x_0,x_0,x_1,\ldots,x_{N-1},x_N,x_N)$ and
\begin{equation}
  f_h(x)=\sum_{i=0}^{N}\theta_i B_i(x),
  \qquad \theta_i=f_h(x_i).
\end{equation}
The exact antiderivative $U_h(x)=\int_0^x f_h(s)\,ds$
belongs to the quadratic spline space obtained by adding
one copy of each endpoint knot. Write
\begin{equation}
  U_h(x)=\sum_{j=0}^{N+1}\beta_j\Phi_j(x),
  \qquad \beta=A_h\theta.
\end{equation}
Define weights (valid also for nonuniform cells)
\begin{equation}
  w_i=\begin{cases}
    h_0/2,           & i=0,       \\
    (h_{i-1}+h_i)/2, & 1\leq i<N, \\
    h_{N-1}/2,       & i=N.
  \end{cases}
\end{equation}
Then integration gives the exact coefficient transformation
\begin{equation}
  \boxed{
    \beta_0=0,\quad
    \beta_j=\sum_{i=0}^{j-1}w_i\theta_i,\quad
    (A_h)_{ji}=\begin{cases}w_i,&i<j,\\0,&i\geq j.\end{cases}}
  \label{eq:scho-weighted-prefix}
\end{equation}
Here $0\leq i\leq N$ and $0\leq j\leq N+1$.
Although $A_h$ has a lower-triangular pattern, we need not
construct a dense matrix: a weighted cumulative sum computes
the coefficients in $O(N)$ operations, without training.

For two interior knots $(x_1,x_2)=(1/3,2/3)$:
\begin{equation}
  A_h=
  \begin{pmatrix}
    0       & 0       & 0       & 0       \\
    \frac16 & 0       & 0       & 0       \\
    \frac16 & \frac13 & 0       & 0       \\
    \frac16 & \frac13 & \frac13 & 0       \\
    \frac16 & \frac13 & \frac13 & \frac16
  \end{pmatrix},\qquad
  A_h\begin{pmatrix}0\\3\\0\\0\end{pmatrix}
  =\begin{pmatrix}0\\0\\1\\1\\1\end{pmatrix}.
  \label{eq:scho-prefix-example}
\end{equation}
The five output numbers are \emph{quadratic spline
  coefficients}, not nodal samples. This reconstructs the
antiderivative exactly; its value at $x=1$ is $1$.

\subsubsection{What about the tridiagonal Poisson matrix?}

For the Dirichlet problem $-u''=f$ on $[0,1]$, the
same degree-one hats $\phi_1,\ldots,\phi_{N-1}$ yield
the Galerkin stiffness matrix
\begin{equation}
  K_{ij}=\int_0^1\phi_i'\phi_j'\,dx,\qquad
  K_{ii}=\frac1{h_{i-1}}+\frac1{h_i},\qquad
  K_{i,i+1}=-\frac1{h_i},
  \label{eq:scho-p1-stiffness}
\end{equation}
with $K_{ij}=0$ when $|i-j|>1$, and symmetric
off-diagonal entries. Hence $K$ is tridiagonal and
$Kc=b$ can be solved in $O(N)$ time. The stiffness
matrix $K$ is \emph{not} the integration map $A_h$:
they serve different purposes. Higher-degree B-splines
keep local support but produce wider bands; they need
not be computationally preferable when linear hats
already achieve the desired accuracy.


\subsection{An interpolation defect and a learned integral correction}
\label{sec:scho-interpolation-defect}

The earlier orthogonal SCHO-Res decomposition concerns
the \emph{output} space. For numerical integration
there is another useful, but non-orthogonal, decomposition
of the \emph{input}. Given an unknown continuous $f$,
let $P_hf$ be its nodal piecewise-linear interpolant:
\begin{equation}
  \boxed{f=P_hf+\varepsilon_f,\quad
    \varepsilon_f=f-P_hf,\quad \varepsilon_f(x_i)=0.}
\end{equation}
The fixed interpolation operator makes this split unique.
It is not an orthogonal projection. By linearity,
\begin{equation}
  \boxed{
    \int_0^x f(s)\,ds=
    \underbrace{\Phi(x)^{\mathsf T}A_h\theta}_{\text{exact base}}
    +\underbrace{\int_0^x\varepsilon_f(s)\,ds}_{\text{unknown residual}},
    \qquad\theta_i=f(x_i).
  }
  \label{eq:scho-split-integral}
\end{equation}
The node values alone do \emph{not} determine the residual
integral for an arbitrary continuous $f$.

\subsubsection{Local reference-cell representation}

For a cell $[x_i,x_{i+1}]$, set
\begin{equation}
  t=(x-x_i)/h_i,\quad g_i(t)=f(x_i+h_it),\quad
  \ell_i(t)=(1-t)\theta_i+t\theta_{i+1}.
\end{equation}
The local defect is $e_i(t)=g_i(t)-\ell_i(t)$.
Instead of learning $e_i$ and integrating a network at
inference, learn its \emph{accumulated} correction:
\begin{equation}
  E_i(t)=\int_0^t e_i(s)\,ds,\qquad
  \boxed{E_i(0)=0,\quad E_i'(t)=e_i(t).}
  \label{eq:scho-defect-ode}
\end{equation}
Since $e_i(0)=e_i(1)=0$, its primitive has
$E_i'(0)=E_i'(1)=0$. Importantly, $E_i(1)$
may be nonzero: it is the signed cell integral error.

An example with just two fixed correction modes is
\begin{equation}
  \psi_0(t)=3t^2-2t^3,\qquad
  \psi_1(t)=t^2(1-t)^2.
\end{equation}
Both vanish at $t=0$ and have zero derivative at
the two endpoints, but $\psi_0(1)=1$.
Given observed cell data, for example
\begin{equation}
  y_i=\bigl(g_i(0),g_i(\tfrac12),g_i(1),h_i\bigr),
\end{equation}
let a small hypernetwork output correction coefficients:
\begin{equation}
  \boxed{\gamma_i=H_\alpha(y_i),\qquad
    \widehat E_i(t)=\sum_{k=0}^{r-1}\gamma_{i,k}\psi_k(t)
    =\Psi(t)^{\mathsf T}H_\alpha(y_i).}
  \label{eq:scho-cell-hypernetwork}
\end{equation}
This is an operator from a finite sensor vector to
a whole local correction function.

\subsubsection{Derivative loss and gradient through the network}

During training, evaluate the true $g_i(t)$ at additional
collocation points. No exact \emph{integral labels} are
needed to enforce the differential definition:
\begin{equation}
  \boxed{\mathcal L(\alpha)=\mathbb E_{i,t}
  \left|\Psi'(t)^{\mathsf T}H_\alpha(y_i)
  -[g_i(t)-\ell_i(t)]\right|^2.}
  \label{eq:scho-defect-loss}
\end{equation}
The network need not differentiate through spline knots
to get the spatial derivative: the fixed $\psi_k'(t)$
are known analytically. Automatic differentiation
does backpropagate through the \emph{coefficient map}.
For the scalar error
$\rho=\Psi'(t)^{\mathsf T}H_\alpha(y_i)-e_i(t)$,
\begin{equation}
  \boxed{\nabla_\alpha(\rho^2)=2\rho\,
    J_{H_\alpha,\alpha}(y_i)^{\mathsf T}\Psi'(t).}
  \label{eq:scho-defect-backprop}
\end{equation}
Thus the path is: loss, correction coefficients,
hypernetwork weights. Here the Jacobian is
$r\times\dim(\alpha)$. Finite sensors do not make an
arbitrary input function identifiable; the training
collocation values supply extra information.

For $x=x_j+h_jt$ the global prediction becomes
\begin{equation}
  \boxed{\widehat U(x)=\Phi(x)^{\mathsf T}A_h\theta
    +\sum_{i<j}h_i\widehat E_i(1)+h_j\widehat E_j(t).}
  \label{eq:scho-defect-global}
\end{equation}
The factor $h_i$ comes from the affine substitution.
No numerical integration of the network is needed
during inference.

\paragraph{A crucial non-neural baseline.}
If $e_i(t)=c_i\,t(1-t)$ and $z_i=g_i(1/2)$, then
$c_i=4[z_i-(\theta_i+\theta_{i+1})/2]$ and
$E_i(t)=(c_i/6)\psi_0(t)$ exactly. Therefore
\begin{equation}
  \int_{x_i}^{x_{i+1}}f(x)\,dx
  =\frac{h_i}{6}(\theta_i+4z_i+\theta_{i+1}),
  \label{eq:scho-simpson}
\end{equation}
which is precisely \emph{Simpson's rule}. This family
does not justify a neural network: with three sensors,
the classical rule is already exact. For other function
families, the speed, extra sampling cost, accuracy and
failure cases must be benchmarked against classical
quadrature. These are proposed training equations,
not experimental results.




\section{PDE-Based Learning without Solution Labels}
\label{sec:scho-physics}

Consider a family of Dirichlet Poisson problems:
\begin{equation}
  -u''(x)=f_\theta(x)\ \ (0<x<1),\qquad u(0)=u(1)=0.
  \label{eq:scho-poisson}
\end{equation}
Its weak form is
\begin{equation}
  \int_0^1u'(x)v'(x)\,dx=\int_0^1 f_\theta(x)v(x)\,dx
  \quad\text{for all }v\in H_0^1(0,1).
  \label{eq:scho-weak}
\end{equation}
Choose a finite set of test functions $v_1,\ldots,v_L$ and a
distribution of input coefficients $\theta$. Define a physics-informed
loss using the predicted spline plus correction:
\begin{equation}
  \boxed{\mathcal L(\alpha,\varphi)=
    \mathbb E_\theta\sum_{\ell=1}^{L}
    \left|\int_0^1\widehat u_\theta'v_\ell'\,dx
    -\int_0^1f_\theta v_\ell\,dx\right|^2.}
  \label{eq:scho-weak-loss}
\end{equation}
It uses \emph{no reference solution values}, but it still requires
the PDE, boundary conditions, forcing samples or coefficients,
test functions, and numerical integration. A low loss over finite
test functions does not certify continuum solution accuracy.

\subsection{Differentiating through the operator}

Writing $\Phi=(\Phi_1,\ldots,\Phi_n)^\mathsf T$ and
$\Psi=(\psi_1,\ldots,\psi_r)^\mathsf T$, the sensitivity to input
coefficients is
\begin{equation}
  \frac{\partial\widehat u_\theta(x)}{\partial\theta}
  =\Phi(x)^\mathsf T J_{H_\alpha}(\theta)
  +\Psi(x)^\mathsf T J_{D_\varphi}(\theta).
  \label{eq:scho-jacobian}
\end{equation}
Automatic differentiation can backpropagate the weak-form loss through
both networks, while spatial derivatives of fixed B-splines can be
evaluated analytically. Quadrature and projection operations must
preserve the required differentiation paths.




\section{Initial Demonstrations and Comparisons}
\label{sec:scho-roadmap}

\begin{enumerate}
  \item Implement a fixed-knot B-spline basis, its derivatives and
        Dirichlet constraints; test coordinate reconstruction and basis rank.
  \item Verify the integral example against the exact matrix $A$.
  \item Train SCHO-Base on the Poisson family using the weak-form loss;
        independently check errors against a numerical spline/Galerkin solver.
  \item Add a few orthogonally projected correction modes, first with
        $H_\alpha$ frozen, then optionally with joint training.
        Compare against simply enlarging the ordinary spline space.
  \item Study a nonlinear coefficient map using, for instance,
        $-u''+u^3=f_\theta$ with homogeneous Dirichlet conditions.
\end{enumerate}
Compare final solution error, distribution shifts, total compute and
memory, training cost and inference latency. A strong classical
discretization is a required baseline: for linear Poisson equations
the discrete coefficient map is itself linear. No improvement,
theorem of universal approximation, or numerical result is claimed
at this specification stage.
